\section{数据与任务噪声：污染、难度漂移与分布错配}

\subsection{数据污染（contamination）}
模型训练语料与 benchmark 测试集重叠，会显著高估泛化能力。课程中给出的核心观点是：污染检测不应只做 exact match，还应关注 paraphrase 和近重复。

\begin{importantbox}{污染治理的三层方案}
\begin{enumerate}
    \item \textbf{静态去重}：n-gram / MinHash / embedding 近邻过滤。
    \item \textbf{过程审计}：训练数据版本化与可追踪清单。
    \item \textbf{评测防御}：动态更新测试集或构造 holdout private split。
\end{enumerate}
\end{importantbox}

\subsection{难度漂移（difficulty drift）}
随着模型能力提升，旧 benchmark 可能迅速饱和。此时分数增益压缩，噪声占比上升，导致 ``谁第一'' 变得更多取决于方差而不是能力差。Sida 建议对 benchmark 进行 difficulty stratification（按难度分桶）并持续更新 hard split。

\begin{knowledgebox}{难度分桶的收益}
分桶后可以区分两类改进：
\begin{itemize}
    \item 在 easy split 提升：可能只是模板适配或表层优化。
    \item 在 hard split 提升：更可能代表真实能力增长。
\end{itemize}
\end{knowledgebox}

\subsection{任务分布错配}
公开 benchmark 的任务分布往往与真实应用分布不同。对于 agent 系统，这种错配更严重，因为线上任务包含权限约束、工具失败、长上下文污染和多轮中断等现实因素。课程强调，benchmarks 应和 production telemetry 联动，而不是孤立存在。

\begin{warningbox}{只优化公开基准会产生局部最优}
模型可能在 benchmark 上显著提升，但在真实任务链路里因为异常恢复和工具鲁棒性不足而退化。
\end{warningbox}

\subsection{本章小结}
数据与任务层噪声决定了 benchmark ``测到的到底是不是你想测的能力''。去污染、难度分桶和分布对齐是提升评测信度的三条主线。

